%\documentclass[a4paper] %,UKenglish,cleveref, autoref, thm-restate]

\documentclass[11pt, oneside]{article}   	
\usepackage{geometry,amsmath,amssymb,amsthm}                		
\geometry{letterpaper}  

\usepackage{tikz}
%This is a template for producing LIPIcs articles. 
%for A4 paper format use option "a4paper", for US-letter use option "letterpaper"
%for british hyphenation rules use option "UKenglish", for american hyphenation rules use option "USenglish"
%for section-numbered lemmas etc., use "numberwithinsect"
%for enabling cleveref support, use "cleveref"
%for enabling autoref support, use "autoref"
%for anonymousing the authors (e.g. for double-blind review), add "anonymous"
%for enabling thm-restate support, use "thm-restate"
%for enabling a two-column layout for the author/affilation part (only applicable for > 6 authors), use "authorcolumns"
%for producing a PDF according the PDF/A standard, add "pdfa"

%\pdfoutput=1 %uncomment to ensure pdflatex processing (mandatatory e.g. to submit to arXiv)
%\hideLIPIcs  %uncomment to remove references to LIPIcs series (logo, DOI, ...), e.g. when preparing a pre-final version to be uploaded to arXiv or another public repository

%\graphicspath{{./graphics/}}%helpful if your graphic files are in another directory

\bibliographystyle{plainurl}% the mandatory bibstyle
\newenvironment{toappendix}{}{}

\newcommand{\bcount}{\mathsf{bcount}}




\newcommand{\longv}[1]{}
\newcommand{\new}[1]{\textcolor{blue}{#1}}

%%% PROVISIONAL COMMANDS
\newcommand{\fL}{\lambda}
\newcommand{\loglog}{\ell_2}
\newcommand{\ACDL}{\mathbb{ACDL}}




\newcommand{\Eval}{\mathsf{Eval}}
\newcommand{\invEval}{\mathsf{invEval}}
\newcommand{\minF}{\mathsf{min}}
\newcommand{\pmin}{\mathsf{pmin}}

\newcommand{\Maj}{\textsc{Maj}}



\title{Characterizing Small Circuit Classes from $\AC^0$ to $\AC^1$ via Discrete Ordinary Differential Equations} 

\author{Melissa Antonelli$^{1,2}$ \quad Arnaud Durand$^3$ \quad Juha Kontinen$^2$}

\date{\small{$^1$Helsinki Institute for Information Technology, $^2$University of Helsinki, $^3$Universit\'e Paris Cit\'e}}



%\titlerunning{Characterizations of Small Circuit Classes via Discrete Ordinary Differential Equations}


%\author{Anonymous Authors}{Anonymous Institution}{}{}{}
%\authorrunning{Anonymous Authors}


%\longv{
%\author{Melissa Antonelli %{Open Access}
%}{HIIT $\&$ University of Helsinki, Pietari Kalmin katu 5, Helsinki Finland}{melissa.antonelli@helsinki.fi}{https://orcid.org/...}{Supported by the Helsinki Institute of Information Technology (HIIT) 
%Helsinki Institute for Information Technology (HIIT) 
%and the Magnus Ehrnrooth Foundation} %(Institute Francais de Finlande, Embassy of France in Finland, French Ministry of Higher Education and Research and Finnish Society of Sciences and Letters).}



%\author{Arnaud Durand %{Open Access}
%}{Université Paris Cité, 8, place Aur\'elie Nemours, Paris, France}{durand@math.univ-paris-diderot.fr}{https://orcid.org/...}{Supported by the Magnus Ehrnrooth Foundation}%TODO mandatory, please use full name; only 1 author per \author macro; first two parameters are mandatory, other parameters can be empty. Please provide at least the name of the affiliation and the country. The full address is optional. Use additional curly braces to indicate the correct name splitting when the last name consists of multiple name parts.



%\author{Juha Kontinen%{Open Access}
%}{University of Helsinki, Pietari Kalmin katu 5, Helsinki, Finland}{juha.kontinen@helsinki.fi}{https://orcid.org/...}{Supported by  the Magnus Ehrnrooth Foundation} %please use full name; only 1 author per \author macro; first two parameters are mandatory, other parameters can be empty. Please provide at least the name of the affiliation and the country. The full address is optional. Use additional curly braces to indicate the correct name splitting when the last name consists of multiple name parts.



%\authorrunning{M. Antonelli, A. Durand and J. Kontinen} %TODO mandatory. First: Use abbreviated first/middle names. Second (only in severe cases): Use first author plus 'et al.'



%\Copyright{CC-BY} %TODO mandatory, please use full first names. LIPIcs license is "CC-BY";  http://creativecommons.org/licenses/by/3.0/
%
%
%\ccsdesc[100]{Theory of computation $\to$ Complexity classes; Theory of computation $\to$ Circuit complexity} %TODO mandatory: Please choose ACM 2012 classifications from https://dl.acm.org/ccs/ccs_flat.cfm 

% end of \longv }


%\keywords{Implicit computational complexity, parallel computation, function algebras, ordinary differential equations, circuit complexity}



%\category{} %optional, e.g. invited paper
%
%\relatedversion{} %optional, e.g. full version hosted on arXiv, HAL, or other respository/website
%\relatedversiondetails[linktext={opt. text shown instead of the URL}, cite=DBLP:books/mk/GrayR93]{Classification (e.g. Full Version, Extended Version, Previous Version}{URL to related version} %linktext and cite are optional
%
%\supplement{}%optional, e.g. related research data, source code, ... hosted on a repository like zenodo, figshare, GitHub, ...
%\supplementdetails[linktext={opt. text shown instead of the URL}, cite=DBLP:books/mk/GrayR93, subcategory={Description, Subcategory}, swhid={Software Heritage Identifier}]{General Classification (e.g. Software, Dataset, Model, ...)}{URL to related version} %linktext, cite, and subcategory are optional
%
%\funding{(Optional) general funding statement \dots}%optional, to capture a funding statement, which applies to all authors. Please enter author specific funding statements as fifth argument of the \author macro.
%
%
%\acknowledgements{Melissa Antonelli and Arnaud Durand thank Maupertuis' Programm for support provided by the Institut Francais de Finlande, the Embassy of France in Finland, the French Ministry of Higher Education and Research, and the Finnish Society of Sciences and Letters for finantially supporting their research\dots}%optional
%
%\nolinenumbers %uncomment to disable line numbering




%\longv{
%Editor-only macros:: begin (do not touch as author)%%%%%%%%%%%%%%%%%%%%%%%%%%%%%%%%%%
%\EventEditors{editors}
%\EventNoEds{...}
%\EventLongTitle{50th International Symposium on Mathematical Foundations of Computer Science (MFCS 2025)}
%\EventShortTitle{MFCS 2025}
%\EventAcronym{MFCS}
%\EventYear{2025}
%\EventDate{..., 2025}
%\EventLocation{...}
%\EventLogo{}
%\SeriesVolume{...}
%\ArticleNo{...}
%}






\usepackage{amsthm,xcolor,hyperref}




\newcommand{\mel}[1]{\marginpar{\textcolor{blue}{Melissa: #1}}}
\newcommand{\arn}[1]{\marginpar{\textcolor{green}{A: #1}}}
\newcommand{\juha}[1]{\marginpar{\textcolor{purple}{J: #1}}}


\newcommand{\wkODE}{wk\text{-ODE}_3}
\newcommand{\ODEfour}{\ell\text{-ODE}_4}
\newcommand{\uptimes}{\times_{\uparrow}}

\newcommand{\ODEzero}{\ell\text{-ODE}_0}
    
\newcommand{\bexp}{\mathsf{bexp}}
\newcommand{\BIT}{\mathsf{BIT}}
\newcommand{\CRN}{\text{CRN}}
\newcommand{\fun}[1]{\mathsf{#1}}








\newcommand{\sB}{B}
\newcommand{\sA}{A}

\newcommand{\Nat}{\mathbb{N}}

%%%%%% CLASSES
\newcommand{\FP}{\mathbf{FP}}
\newcommand{\NC}{\mathbf{FNC}}
\newcommand{\AC}{\mathbf{FAC}}
\newcommand{\TC}{\mathbf{FTC}}
\newcommand{\FACC}{\mathbf{FACC[2]}}

%%%%%% CIRCUITS
\newcommand{\size}{\textsf{size}}
\newcommand{\depth}{\textsf{depth}}



%%% SIGN / TERMS
\newcommand{\sign}{\mathsf{sg}}
\newcommand{\term}{\textsf{TERM}}
\newcommand{\degre}{\textsf{deg}}

\newcommand{\termO}{s}
\newcommand{\termT}{t}
\newcommand{\xO}{\mathbf{x}}
\newcommand{\yO}{\mathbf{y}}

\newcommand{\signTerm}{\textsf{SIGN}\text{-}\textsf{TERM}}

\newcommand{\expr}{T}
\newcommand{\expO}{A}
\newcommand{\expT}{B}


\newcommand{\fO}{\mathbf{f}}
\newcommand{\fT}{\mathbf{g}}
\newcommand{\fTh}{\mathbf{h}}
\newcommand{\fF}{\mathbf{u}}

\newcommand{\Bm}{\mathbf{B}}



%%%% ODE
\newcommand{\llLODE}{\ell^2\text{-}lODE}
\newcommand{\lLODE}{\ell\text{-}lODE}
\newcommand{\bcODE}{\ell^2\text{-}bcODE}

\newcommand{\LDL}{\mathbb{LDL}}

\newcommand{\Zf}{\mathbf{0}}
\newcommand{\Of}{\mathbf{1}}

\newcommand{\tu}[1]{\mathbf{#1}}


%%%% MAJ TH
\newcommand{\maj}{\textsc{Maj}}
\newcommand{\thGate}{\textsc{Th}}
\newcommand{\bit}{\textsf{bit}}
\newcommand{\modTwo}{\textsf{mod 2}}
\newcommand{\sucZero}{\mathsf{s}_0}
\newcommand{\sucOne}{\mathsf{s}_1}
\newcommand{\sucI}{\mathsf{s}_i}
\newcommand{\zeroF}{\mathsf{0}}
\newcommand{\oneF}{\mathsf{1}}
\newcommand{\MOD}{\textsf{ mod 2}}
\newcommand{\msp}{\textsf{msp}}
\newcommand{\MSP}{\textsf{MSP}}
\newcommand{\IF}{\textsf{if}}
\newcommand{\cond}{\textsf{cond}}


%%% ODE
\newcommand{\ODEone}{\ell\text{-ODE}_1}
\newcommand{\ODEtwo}{\ell\text{-ODE}_2}
\newcommand{\ODEthree}{\ell\text{-ODE}_3}

\newcommand{\TCDL}{\mathbb{TCDL}}
\newcommand{\TClass}{\mathbf{FTC}}


\newcommand{\circACDL}{\ACDL_{\mathcal{C}}}
\newcommand{\circTCDL}{\TCDL_{\mathcal{C}}}
\newcommand{\circuitfunction}{\mathbf{circ}_{\mathcal{C}}}


%%%% NEWTHEOREM
\newtheorem{prop}{Proposition}
\newtheorem{theorem}{Theorem}
\newtheorem{lemma}{Lemma}

\theoremstyle{definition}
\newtheorem{defn}{Definition}
\newtheorem{notation}{Notation}
\newtheorem{example}{Example}
\newtheorem{remark}{Remark}









%%%
\newcommand{\ODEACC}{\ell\text{-ODE}_{\text{FACC[2]}}}





\newcommand{\schODE}{\ell\text{-b}_0\text{ODE}}


\newcommand{\chODE}{\ell\text{-bODE}}

\newcommand{\pODE}{\ell\text{-pODE}}





































\usepackage{fontspec}
